\begin{figure*}[t]
\centering
\scriptsize
\setlength{\fboxsep}{1em}
\setlength{\fboxrule}{1pt}
\fcolorbox{black}{white}{
\parbox{0.955\textwidth}{
\begin{alltt}


You will be provided instructions for completing a task followed by a task to complete.

Instructions:

1. Understand the Passage: The passage will either be a brief summary of the plot or the full story. Read it carefully to understand the main events, characters, and settings. 

2. Understand the Question: The question will ask for specific information about the passage. It could be about a character, an event, or a detail about the setting. Make sure you understand what the question is asking.

3. Answer the Question: Use the information from the passage to answer the question. Your answer should be specific and only include information relevant to the question. Do not include additional information not asked for in the question.

4. Be Concise: The answer should be concise and to the point. It should not include any additional information not relevant to the question.

5. Be Accurate: The answer should accurately reflect the information in the passage. Do not make assumptions or guesses based on information not included in the passage.

6. Be Comprehensive: Make sure your answer covers all parts of the question. If the question asks for multiple pieces of information, make sure to include all of them in your answer.

7. Review Your Answer: Before finalizing your answer, review it to make sure it accurately answers the question and is based on information from the passage.

Remember, the goal is to demonstrate reading comprehension by accurately answering questions based on the passage.

Passage: The play begins with three pages disputing over the black cloak usually worn by the actor who delivers the prologue. They draw lots for the cloak, and one of the losers, Anaides, starts telling the audience what happens in the play to come; the others try to suppress him, interrupting him and putting their hands over his mouth. Soon they are fighting over the cloak and criticizing the author and the spectators as well.

In the play proper, the goddess Diana, also called Cynthia, has ordained a "solemn revels" in the valley of Gargaphie in Greece. The gods Cupid and Mercury appear, and they too start to argue. Mercury has awakened Echo, who weeps for Narcissus, and states that a drink from Narcissus's spring causes the drinkers to "Grow dotingly enamored of themselves." The courtiers and ladies assembled for the Cynthia's revels all drink from the spring.

Asotus, a foolish spendthrift who longs to become a courtier and a master of fashion and manners, also drinks from the spring; emboldened by vanity and self-love, he challenges all comers to a competition of "court compliment." The competition is held, in four phases, and the courtiers are beaten. Two symbolic masques are performed within the play for the assembled revelers. At their conclusion, Cynthia (representing Queen Elizabeth) has the dancers unmask and shows that vices have masqueraded as virtues. She sentences them to make reparation and to purify themselves by bathing in the spring at Mount Helicon.

The figure of Actaeon in the play may represent Robert Devereux, 2nd Earl of Essex, while Cynthia's lady in waiting Arete may be Lucy, Countess of Bedford, one of Elizabeth's ladies in waiting as well as Jonson's patroness.

The play is notably rich in music, as is typical for the theatre of the boys' companies, which originated as church choirs.

Question: WHAT NAME WAS CYNTHIA MORE FAMOUSLY KNOWN BY?

Follow the instructions to generate an explanation that reasons towards the correct answer to the task above. End the explanation with the correct answer. [/INST] Explanation: 

The passage describes the play "The Masque of Queens" by Ben Jonson, which features the goddess Diana, also known as Cynthia, who ordains a "solemn revels" in the valley of Gargaphie in Greece. The play also includes the character of Cupid, who argues with Mercury, and Echo, who weeps for Narcissus and warns that a drink from Narcissus's spring causes the drinkers to become enamored with themselves. The play features a competition of "court compliment" in which the foolish spendthrift Asotus challenges all comers and is beaten, and includes two symbolic masques performed for the assembled revelers. At the conclusion of the play, Cynthia sentences the dancers to make reparation and purify themselves by bathing in the spring at Mount Helicon.

The question asks for the name that Cynthia is more famously known by. Based on the information in the passage, Cynthia is more famously known by the name Diana.

Answer: Diana.

Therefore, the answer to the task is below. Give the answer in the shortest form possible that will still be correct.

Answer: Diana.
\end{alltt}
}}
\caption{A complete running example of NarrativeQA.}
\label{fig:narrative_qa_example}
\end{figure*}
